\section{Discussion}\label{sec:discussion}
\noindent\textbf{What the network adds.} Every result in Sections~\ref{sec:headtohead} to \ref{sec:hardware} concerns a quantity that does not exist in a dataset evaluation: which legitimate devices are disconnected, whether the detector keeps working when its own queue fills, how long telemetry waits behind a flood, and how much inspection energy an inert detector wastes. Three rules with identical detection rates differ between quarantining none of a ward's life-critical devices and quarantining every device on the ward. A rule that looks acceptable in the baseline switches itself off when the same traffic is routed to fewer fog nodes. A detector that keeps detecting under a realistic tail inference cost still delays wearable telemetry by most of a second. None of this is a criticism of the audited authors' care; it is what the evaluation instrument cannot see.

\noindent\textbf{What the audit adds.} The arithmetic checks need no code and no data, and they caught errors in two of three papers. The most consequential, partitioning after oversampling, coincides with the largest reported scores: on WUSTL-EHMS-2020 our own 20-partition replication moves F1 by 0.29 with the order alone, and three independent pipelines land on either side of that gap according to their order. A reader comparing published IoMT detectors on that dataset is largely comparing evaluation protocols.

\noindent\textbf{Duplicates need a feature space.} A duplicate rate stated without its columns cannot be checked. Our own earlier Edge-IIoTset figure was challenged for exactly this reason, and both it and the conflicting published figure turn out to be correct for their populations. We recommend that papers report duplicate and train/test-twin rates in the feature space the model is trained on, which is the only one that bears on evaluation.

\noindent\textbf{Design guidance for in-path detectors.} Three follow from the results. Do not couple a detection threshold to observed congestion without a floor: a detector in the data path controls the congestion it observes, and the loop can lock in the wrong state. If a rule must adapt to load, couple it to measured capacity, so that overload, which is itself evidence of an attack, tightens rather than loosens it. And state per-device consequences: a rule's quarantine profile by device class is a clinical decision that the criticality weights make explicit. At its default settings C never disconnected a life-critical device in any scenario, but a higher base threshold, flatter weights or a mimicry attack did reach that class, so the guarantee belongs to a configuration, not to the rule.

\noindent\textbf{Negative results.} Our ensemble did not beat its best branch, LIME was unstable, the detector did not recognise withheld attack types, a threshold rule alone did not reach its false-positive target in the detection pipeline, and cross-dataset transfer between the requested datasets is impossible without re-extraction. We report these because the literature's ceiling effect is partly made of results like them left unreported.
